% 付録C：基礎ベースライン（求人タイトルのみ・業務内容のみ）

本付録では、複数項目を統合した推薦手法（本文の本実験）に先立つ基礎検証として、\textbf{求人タイトルのみ}および\textbf{業務内容のみ}から応募履歴を学習した場合の推定精度を示す。

\section{求人タイトルのみ}
求人タイトル（職種名）のみを用い、過去応募履歴で出現したタイトルの出現頻度に基づいて候補求人を順位付けした。

\begin{table}[H]
\centering
\caption{ワーカー別の推薦精度（求人タイトルのみ、30件学習、10件テスト）}
\label{tab:appc-worker-accuracy-title-only}
\scriptsize
\begin{tabular}{r l r r r r}
\hline
順位 & ワーカー & Top-1 & Top-3 & Top-5 & 複合スコア \\
\hline
1  & ワーカー A & 100.0\% & 100.0\% & 100.0\% & 1.0000 \\
2  & ワーカー B & 100.0\% & 100.0\% & 100.0\% & 1.0000 \\
3  & ワーカー C & 60.0\%  & 80.0\%  & 80.0\%  & 0.7000 \\
4  & ワーカー D & 50.0\%  & 80.0\%  & 80.0\%  & 0.6500 \\
5  & ワーカー E & 60.0\%  & 60.0\%  & 60.0\%  & 0.6000 \\
6  & ワーカー F & 20.0\%  & 80.0\%  & 80.0\%  & 0.5000 \\
7  & ワーカー G & 20.0\%  & 30.0\%  & 40.0\%  & 0.2700 \\
8  & ワーカー H & 10.0\%  & 30.0\%  & 50.0\%  & 0.2400 \\
9  & ワーカー I & 10.0\%  & 20.0\%  & 30.0\%  & 0.1700 \\
10 & ワーカー J & 0.0\%   & 30.0\%  & 30.0\%  & 0.1500 \\
\hline
全体平均（10名） & --- & 43.0\% & 61.0\% & 65.0\% & 0.5280 \\
\hline
\end{tabular}

\end{table}

\begin{table}[H]
\centering
\caption{学習件数別の精度推移（求人タイトルのみ、20名）}
\label{tab:appc-learning-curve-title-only-20workers}
\scriptsize
\begin{tabular}{r r r r r}
\hline
学習件数 & n & 平均Top-1 & 平均Top-3 & 平均Top-5 \\
\hline
10 & 20 & 0.4500 & 0.5000 & 0.5700 \\
15 & 20 & 0.4600 & 0.6650 & 0.7050 \\
20 & 17 & 0.6059 & 0.7765 & 0.8000 \\
25 & 11 & 0.4909 & 0.6727 & 0.7273 \\
30 & 10 & 0.4300 & 0.6100 & 0.6500 \\
35 & 7 & 0.5429 & 0.7143 & 0.7714 \\
40 & 5 & 0.5400 & 0.7000 & 0.7600 \\
\hline
\end{tabular}

\end{table}

\begin{table}[H]
\centering
\caption{就業回数上位10名を除いた場合の精度推移（求人タイトルのみ）}
\label{tab:appc-learning-curve-title-only-excluding-top10}
\scriptsize
\begin{tabular}{r r r r r}
\hline
学習件数 & n & 平均Top-1 & 平均Top-3 & 平均Top-5 \\
\hline
10 & 10 & 0.5200 & 0.5600 & 0.6500 \\
15 & 10 & 0.5000 & 0.7100 & 0.7300 \\
20 & 7 & 0.7143 & 0.8143 & 0.8143 \\
25 & 1 & 0.1000 & 0.4000 & 0.4000 \\
\hline
\end{tabular}

\end{table}

\section{業務内容のみ}
業務内容テキストのみを用い、内容の近さに基づいて候補求人を順位付けした。

\begin{table}[H]
\centering
\caption{学習件数別の精度推移（業務内容のみ、20名）}
\label{tab:appc-learning-curve-content-only-20workers}
\scriptsize
\begin{tabular}{r r r r r}
\hline
学習件数 & n & 平均Top-1 & 平均Top-3 & 平均Top-5 \\
\hline
10 & 20 & 0.4350 & 0.4950 & 0.5650 \\
15 & 20 & 0.3700 & 0.4450 & 0.5450 \\
20 & 17 & 0.5412 & 0.6235 & 0.7294 \\
25 & 11 & 0.4273 & 0.5182 & 0.6455 \\
30 & 10 & 0.3000 & 0.4000 & 0.5200 \\
35 & 7 & 0.2857 & 0.5286 & 0.5286 \\
40 & 5 & 0.4800 & 0.5400 & 0.6000 \\
\hline
\end{tabular}

\end{table}

\begin{table}[H]
\centering
\caption{就業回数上位10名を除いた場合の精度推移（業務内容のみ）}
\label{tab:appc-learning-curve-content-only-excluding-top10}
\scriptsize
\begin{tabular}{r r r r r}
\hline
学習件数 & n & 平均Top-1 & 平均Top-3 & 平均Top-5 \\
\hline
10 & 10 & 0.4300 & 0.4700 & 0.5600 \\
15 & 10 & 0.3300 & 0.3700 & 0.4800 \\
20 & 7 & 0.6571 & 0.7000 & 0.7286 \\
25 & 1 & 0.1000 & 0.6000 & 0.7000 \\
\hline
\end{tabular}

\end{table}
